% ============================================================================
% Section 3.2 — Adding Integers
% CCSS 7.NS.A.1.B
% ============================================================================
\section{Adding Integers}

\begin{stepGoal}
\begin{itemize}
    \item Add integers using the sign rules
    \item Understand that a number plus its opposite equals zero
\end{itemize}
\end{stepGoal}

\begin{stepVocabBox}
    \vocabItem{Additive Inverse}{A number and its opposite — they always add to $0$.}
    \vocabItem{Absolute Value}{The distance from $0$, used to decide the size of the sum.}
\end{stepVocabBox}

\begin{stepsCard}[How to Add Two Integers]
    \stepItem{Check the \textbf{signs}. Are they the same or different?}
    \stepItem{\textbf{Same signs:} Add the absolute values. Keep the sign.}
    \stepItem{\textbf{Different signs:} Subtract the smaller absolute value from the larger. Keep the sign of the number with the larger absolute value.}
\end{stepsCard}

% --- Worked Example 1 ---
\begin{stepExample}{Find $(-8) + (-5)$.}
    \stepShow{1} Both numbers are \textbf{negative} — same signs.

    \stepShow{2} Add the absolute values: $8 + 5 = 13$. Keep the negative sign.

    \stepShow{3} (Not needed — signs are the same.)
    \stepResult{$(-8) + (-5) = -13$}
\end{stepExample}

% --- Worked Example 2 ---
\begin{stepExample}{Find $(-12) + 7$.}
    \stepShow{1} $-12$ is negative, $7$ is positive — \textbf{different signs}.

    \stepShow{2} (Not needed — signs are different.)

    \stepShow{3} Subtract: $12 - 7 = 5$. The larger absolute value is $12$ (negative), so keep the negative sign.
    \stepResult{$(-12) + 7 = -5$}
\end{stepExample}

\watchOut{Don't just add the numbers and slap on a negative sign. When the signs are different, you \textbf{subtract} the absolute values!}

% ============================================================================
% PRACTICE
% ============================================================================
\newpage
\begin{stepPractice}{Adding Integers Practice}
    \resetProblems

    \stepPracticeHeader[step1Color]{Same or Different Signs?}
    \begin{multicols}{2}
    \prob $(-3) + (-9)$: same or different? \answerBlank[2cm]
    \answer{Same}
    \prob $15 + (-6)$: same or different? \answerBlank[2cm]
    \answer{Different}
    \end{multicols}

    \stepPracticeHeader[step2Color]{Add — Same Signs}
    \begin{multicols}{2}
    \prob $(-4) + (-7) =$ \answerBlank[2cm]
    \answer{$-11$}
    \prob $(-15) + (-3) =$ \answerBlank[2cm]
    \answer{$-18$}
    \end{multicols}

    \stepPracticeHeader[step3Color]{Add — Different Signs}
    \begin{multicols}{2}
    \prob $(-9) + 14 =$ \answerBlank[2cm]
    \answer{$5$}
    \prob $11 + (-11) =$ \answerBlank[2cm]
    \answer{$0$}
    \end{multicols}

    \wordProblem{The temperature is $-6°$C in the morning. By noon it rises $14°$. What is the noon temperature?}{°C}
    \answer{$8°$C}
\end{stepPractice}
